\chapter{Introduction and Orientation}
\label{ch:introduction}

\section{Motivation}

Consider a system of linear equations
\begin{equation}\label{eq:system}
  \left\{\;\begin{aligned}
    a_{11}x_1+\dots+a_{1n}x_n &= b_1,\\
                              &\;\;\vdots\\
    a_{m1}x_1+\dots+a_{mn}x_n &= b_m.
  \end{aligned}\right.
\end{equation}
Every student can solve such a system by elimination. But elimination answers
only the question \emph{what are the solutions?}. It does not answer the
questions an examiner actually asks: for how many right-hand sides $b$ does a
solution exist? When is it unique? If the data $a_{ij}$ are perturbed slightly,
does the answer change? What is common to \eqref{eq:system}, to a linear
differential equation, and to a recurrence relation?

Linear algebra is the theory that replaces the computation by a structural
answer. The set of solutions of \eqref{eq:system} with $b=0$ is a vector space;
its dimension is $n-\rank A$; the system with $b\neq0$ is solvable exactly when
$b$ lies in the image of the associated map. Three notions -- space, map, rank --
absorb the whole question.

\begin{intuition}
Linear algebra studies the transformations of space that send straight lines to
straight lines and fix the origin. Such a transformation may stretch, rotate,
shear or flatten space, but it never bends it. Because the deformation is
uniform, the whole transformation is determined by what it does to finitely many
vectors: that is why the subject is computable, and also why it is everywhere.
\end{intuition}

\section{Historical background}

Determinants appear before matrices: Leibniz uses them in 1693, Cramer publishes
his rule in 1750. The matrix as an object in its own right is due to Cayley
(1858), who also states the theorem that bears his name with Hamilton. The
axiomatic notion of a vector space, freed from coordinates, is much later: it is
made explicit by Peano (1888) and becomes standard only with the treatises of
the 1930s, notably van der Waerden and Banach for the infinite-dimensional case.
The modern teaching order -- spaces first, matrices as a representation -- thus
reverses the historical order. Keep this in mind: coordinates were the beginning
of the subject, not its content.

\section{How to use this book}

The volume follows the fixed structure of the series.

\begin{itemize}
  \item \Cref{ch:definitions} contains the core definitions with their examples
        and, just as importantly, their non-examples.
  \item \Cref{ch:theorems} contains the fundamental theorems with complete proofs.
  \item \Cref{ch:examples} works out detailed computations, with figures.
  \item \Cref{ch:summary} gathers the common mistakes, a summary, and the
        practice exercises.
  \item \Cref{ch:problems} contains doctoral entrance problems, and
        \cref{ch:solutions} their fully worked solutions, together with the
        solutions of the exercises.
\end{itemize}

A reasonable working method: read a section, close the book, reconstruct the
statements from memory, and only then look at the proofs. Solve the exercises of
\cref{ch:summary} before reading any solution; a solution read is worth roughly
a tenth of a solution found.

\begin{note}[Notation]
Throughout, $\K$ denotes an arbitrary field, $E$ and $F$ are $\K$-vector spaces,
and $u,v$ denote linear maps. $\Mat{n}{\K}$ is the algebra of $n\times n$
matrices over $\K$, $\id$ the identity map, $\transpose{A}$ the transpose of $A$,
and $\spec(u)$ the set of eigenvalues of $u$.
\end{note}

\section{Further reading}

For a coordinate-free treatment close in spirit to this volume, see
\textcite{axler2015}. For the matrix-analytic point of view, \textcite{horn2013}
is the standard reference. Readers preparing Algerian doctoral examinations will
recognise the style of exercises in \textcite{gourdon2009} and
\textcite{grifone2011}.
